\chapter*{Introduction générale}
\addcontentsline{toc}{chapter}{Introduction générale}
\markboth{Introduction générale}{Introduction générale}

Les troubles du sommeil constituent un problème de santé publique dont l'ampleur reste
largement sous-estimée. Le syndrome d'apnées obstructives du sommeil, à lui seul,
concernerait près d'un milliard d'adultes dans le monde, dont une large majorité n'a
jamais été diagnostiquée. Son examen de référence, la polysomnographie, enregistre une
nuit complète de signaux physiologiques ; son interprétation impose ensuite à un
technicien qualifié de classer manuellement un millier de fenêtres de trente secondes,
puis de caractériser les événements respiratoires. Ce goulet d'étranglement humain,
davantage que la capacité d'enregistrement, limite aujourd'hui l'accès au diagnostic.

L'apprentissage profond a apporté une réponse crédible au premier problème. Depuis
2017, plusieurs architectures atteignent sur les bases publiques un niveau de
concordance comparable à celui observé entre deux experts humains. Leur adoption
clinique demeure pourtant marginale, et la raison n'est plus la performance : c'est
l'opacité. Un réseau de neurones profond restitue une classe sans énoncer ce qui l'a
déterminée, alors que le praticien engage sa responsabilité sur le compte rendu qu'il
signe. Sans justification intelligible, une prédiction ne peut être ni validée ni
contestée.

Le présent projet de fin d'études, mené au sein du Centre National de l'Informatique,
répond à cette double exigence de performance et de transparence. Il aboutit à
\hypnoria{}, une plateforme web d'aide à l'interprétation polysomnographique organisée
autour de deux étages d'apprentissage. Le premier classe automatiquement les stades du
sommeil à partir d'un enregistrement brut, en combinant par empilement trois familles
d'architectures — convolutive, récurrente et attentionnelle. Le second exploite
l'hypnogramme ainsi produit pour prédire la classe de sévérité du syndrome d'apnées,
et accompagne chaque prédiction de la décomposition des facteurs qui l'ont déterminée.
Autour de ces modèles, la plateforme fournit la chaîne applicative complète attendue en
contexte clinique : dossier patient, archivage sécurisé, annotation, compte rendu
exportable, collaboration entre praticiens et traçabilité des actions.

Ce mémoire rend compte de l'ensemble de ces travaux. Le
chapitre~\ref{chap:cadre} pose le cadre du projet : l'organisme d'accueil, les notions
médicales indispensables, la problématique et l'étude de l'existant. Le
chapitre~\ref{chap:etat-art} dresse l'état de l'art des méthodes de classification
automatique des stades, des techniques d'apprentissage ensembliste et des approches
d'intelligence artificielle explicable. Le chapitre~\ref{chap:analyse} formalise les
besoins, présente le \emph{Product Backlog} et le découpage du projet en six sprints.
Le chapitre~\ref{chap:conception} détaille la conception du système, de l'architecture
logicielle jusqu'aux architectures de réseaux. Le chapitre~\ref{chap:realisation-ia}
expose la réalisation du pipeline d'apprentissage, la stratégie d'entraînement retenue
sous contrainte matérielle et l'ensemble des résultats expérimentaux. Le
chapitre~\ref{chap:realisation-plateforme} décrit la plateforme clinique et son
déploiement. Une conclusion générale dresse enfin le bilan des objectifs atteints et
ouvre sur les perspectives.
